% Generated draft: requires mathematical and source review.
\documentclass[11pt]{article}
\usepackage[margin=25mm]{geometry}
\usepackage{fontspec}
\setmainfont{Latin Modern Roman}
\usepackage{amsmath,amssymb,mathtools,unicode-math}
\setmathfont{Latin Modern Math}
\usepackage{enumitem,needspace,xurl,hyperref}
\setlength{\parindent}{0pt}
\setlength{\parskip}{5pt}
\setlength{\emergencystretch}{3em}
\newcommand{\N}{\mathbb{N}}
\newcommand{\Z}{\mathbb{Z}}
\newcommand{\Q}{\mathbb{Q}}
\newcommand{\R}{\mathbb{R}}
\newcommand{\C}{\mathbb{C}}
\newcommand{\F}{\mathbb{F}}
\newcommand{\A}{\mathbb{A}}
\newcommand{\PP}{\mathbb{P}}
\newcommand{\E}{\mathbb{E}}
\newcommand{\eps}{\varepsilon}
\newcommand{\dd}{\,\mathrm{d}}
\begin{document}
\section*{UnsolvedMath ID: 800580}
\textbf{Source problem code:} catalogue entry
\subsection{Definitions and mathematical statement}
For a class \(\mathcal J\) of sets of reals, define its Borel conjecture by
\[
\operatorname{BC}(\mathcal J):\quad
\forall X\in\mathcal J\; (X\text{ is countable}).
\]
Here countable means finite or countably infinite.

The supplied excerpt does not define its selection-principle notation. The following gives a conventional reconstruction, subject to verification against the source. For \(X\subseteq\R\), let \(\mathrm B(X)\) consist of countable families \(\mathcal U\) of Borel subsets of the subspace \(X\) such that
\[
\bigcup\mathcal U=X
\quad\text{and}\quad
X\notin\mathcal U.
\]
In this reconstruction, countable families may be finite. Say that \(X\) satisfies \(\mathsf S_{\mathrm{fin}}(\mathrm B,\mathrm B)\) if
\[
\forall (\mathcal U_n)_{n\in\N}\in\mathrm B(X)^{\N}\;
\exists (\mathcal F_n)_{n\in\N}\;
\left[
\left(\forall n\in\N,\ \mathcal F_n\subseteq\mathcal U_n
\text{ is finite}\right)
\land
\bigcup_{n\in\N}\mathcal F_n\in\mathrm B(X)
\right].
\]
Empty finite selections are allowed. The union in the last expression is a family of selected Borel sets.

\textbf{Question.} Is Borel's Conjecture for \(\mathsf S_{\mathrm{fin}}(\mathrm B,\mathrm B)\) consistent? Under the usual interpretation that consistency is with \(\mathrm{ZFC}\), the requested consistency statement is
\[
\operatorname{Con}\!\left(
\mathrm{ZFC}+
\left[\forall X\subseteq\R,\quad
X\models\mathsf S_{\mathrm{fin}}(\mathrm B,\mathrm B)
\ \Longrightarrow\ X\text{ is countable}\right]
\right).
\]
The excerpt does not explicitly specify the background theory or the cover conventions used above; these are reconstruction qualifications requiring review.

The source also states that Borel's Conjecture is consistent for all but three properties in its ``Borel case.'' The excerpt does not identify that collection or the three exceptions, so this contextual assertion is retained as supplied context and does not resolve the question.
\subsection{Short English statement}
The question asks whether it is consistent that every set of reals satisfying the Borel finite-selection property is countable. That property permits finitely many choices from each cover in any sequence of countable Borel covers, with the chosen sets together covering the space. The cover conventions and the interpretation of consistency with \(\mathrm{ZFC}\) are reconstructed and require source verification.
\subsection{Sources}
\begin{itemize}
\item Local UnsolvedMath registry: permanent ID \(800580\), exact source problem code \texttt{catalogue entry}. The supplied record labels the problem \texttt{open}; this is supplied registry metadata, not an independently verified status assessment.
\item Tsaban, \emph{Selection Principles and Special Sets of Reals: Open Problems} (2006), Problem 6. Supplied source URL: \url{https://arxiv.org/abs/math/0609601}. The record reports access on 2026-10-16 and extraction by \texttt{arxiv\_latex}; these provenance details are supplied and have not been independently verified.
\item Supplied attribution/context: \texttt{[pawlikowskireclaw]}. The corresponding bibliographic entry is absent from the supplied record, so this key is preserved without expanding its attribution.
\end{itemize}
\end{document}
